\documentclass[11pt,a4paper]{article}
\usepackage[margin=1.8cm]{geometry}
\usepackage{amsmath,amssymb,tabularx,array,needspace,fontspec,cancel}
\setmainfont{FreeSerif}
\setlength{\parindent}{0pt}
\setlength{\parskip}{0pt}
\renewcommand{\arraystretch}{1.15}
\newcommand{\mcq}[6]{%
  \Needspace{4\baselineskip}%
  \par\medskip\noindent\textbf{#1.}\ #2\par\nopagebreak
  \noindent\begin{tabularx}{\linewidth}{@{\hspace{1.5em}}*{4}{>{\raggedright\arraybackslash}X}@{}}
    (A)~#3 & (B)~#4 & (C)~#5 & (D)~#6
  \end{tabularx}\par}

\begin{document}
\begin{center}
  {\Huge\bfseries The Hope Science Academy}\\[3pt]
  {\small Excellence is our standard not goal}\\[2pt]
  {\small 247-E-1, (Opposite LDA School) Johar Town, Lahore. PH: 0300-8194789, 0322-4157001}
\end{center}
\medskip\noindent
\begin{tabularx}{\linewidth}{|X|X|X|X|}
\hline
\textbf{Student Name:} & \textbf{Roll Num:} & \textbf{Class:} 10th & \textbf{Paper Code:} \\ \hline
\textbf{Subject Name:} Mathematics. & \textbf{Time Allowed:} & \textbf{Total Marks:} 110 & \textbf{Exam Date:} 09-Oct-2026 \\ \hline
\multicolumn{4}{|l|}{\textbf{Exam Syllabus:} FULL BOOK} \\ \hline
\end{tabularx}
\bigskip
\noindent\textbf{Q1. Choose the correct answer.}\hfill\textbf{1$\times$110=110}

\mcq{1}{$i^{2}+i^{4}=$}{-1}{0}{1}{2}
\mcq{2}{Real part of $(2-3i)(2+3i)$}{-3}{1}{4}{13}
\mcq{3}{Imaginary part of $(2-i)(2+i)$ is:}{0}{1}{7}{9}
\mcq{4}{$x+iy$ will be pure imaginary number, when:}{y=0}{x=0}{$i=0$}{x=0, y=0}
\mcq{5}{What is additive inverse of $5-2i$?}{$5+2i$}{$-5-2i$}{$5-2i$}{$-5+2i$}
\mcq{6}{What is multiplicative inverse of $z=1+i$?:}{$1-i$}{$i$}{$\dfrac{1}{2}-\dfrac{1}{2}i$}{$-\dfrac{1}{\sqrt{2}}+\dfrac{1}{\sqrt{2}}i$}
\mcq{7}{If $z=4-3i$, then $z\bar{z}=$}{3}{9}{16}{25}
\mcq{8}{Conjugate of $9-4i$:}{$-9-4i$}{$9+4i$}{$9+9i$}{$4-9i$}
\mcq{9}{If $z=4+4i$, then $z+\bar{z}=$}{8}{$8+8i$}{$8i$}{0}
\mcq{10}{If $z=5+4i$, then $|z|=$}{9}{25}{41}{$\sqrt{41}$}
\mcq{11}{The type of equation $2x^{2}-x+1=0$:}{Quadratic}{Linear}{Third degree}{Pure quadratic}
\mcq{12}{What is the discriminant of $x^{2}+5x-5=0$}{-20}{20}{25}{45}
\mcq{13}{The solution set of $3x^{2}-9=0$:}{\{3\}}{\{$\pm3$\}}{\{$\pm\sqrt{3}$\}}{\{$\sqrt{3}$\}}
\mcq{14}{Sum of the roots of $3x^{2}+5x-12=0$:}{$\dfrac{5}{3}$}{$-\dfrac{5}{3}$}{$\dfrac{12}{3}$}{$\dfrac{3}{5}$}
\mcq{15}{Product of the roots of $3x^{2}+5x-12=0$:}{-4}{3}{4}{5}
\mcq{16}{What are roots of $(x-3)(x+3)=0$:}{3,-3}{3,3}{-3,-3}{9,0}
\mcq{17}{2 and 3 are the roots of:}{$x^{2}+5x+6=0$}{$x^{2}+6x+5=0$}{$x^{2}-5x+6=0$}{$x^{2}+6x-5=0$}
\mcq{18}{If $b^{2}-4ac>0$ and is a perfect square, then the roots of $ax^{2}+bx+c=0$ are:}{Equal}{Unequal}{Imaginary}{Irrational}
\mcq{19}{If $b^{2}-4ac=0$ then the roots of $ax^{2}+bx+c=0$ are:}{Unequal}{Irrational}{Imaginary}{Equal}
\mcq{20}{Subject ``C'' of $x-2c=b$ is:}{x+b}{b-x}{$\dfrac{x-b}{2}$}{$\dfrac{b-x}{2}$}
\mcq{21}{If $\begin{bmatrix}a+2\\3\end{bmatrix}=\begin{bmatrix}7\\3\end{bmatrix}$ then a=:}{3}{5}{6}{7}
\mcq{22}{$A=\begin{bmatrix}3&5&0\end{bmatrix}$ is a ....... matrix:}{Row}{Square}{Column}{Null}
\mcq{23}{$B=\begin{bmatrix}2\\3\end{bmatrix}$ is a ....... matrix:}{Identity}{Square}{Row}{Column}
\mcq{24}{$M=\begin{bmatrix}1&0&0\\0&1&0\\0&0&1\end{bmatrix}$ is a/an ....... matrix:}{Rectangular}{Identity}{Column}{Row}
\mcq{25}{If $A^{t}=-A$, then A is ....... matrix:}{Symmetric}{Row}{Rectangular}{Skew-symmetric}
\mcq{26}{If $A=\begin{bmatrix}3&4\end{bmatrix}$, $B=\begin{bmatrix}7&8\end{bmatrix}$, then $A+B=$}{[ 21 32]}{[24 28]}{[10 12]}{[11 11]}
\mcq{27}{If $A=\begin{bmatrix}3&4\end{bmatrix}$, $B=\begin{bmatrix}7&8\end{bmatrix}$, then $B-A=$}{[10 12]}{[11 11]}{[4 4]}{[-4 -4]}
\mcq{28}{What is the additive of $\begin{bmatrix}3\\4\end{bmatrix}$?}{$\begin{bmatrix}-3\\4\end{bmatrix}$}{$\begin{bmatrix}-3\\-4\end{bmatrix}$}{$\begin{bmatrix}3\\4\end{bmatrix}$}{$\begin{bmatrix}-1\\1\end{bmatrix}$}
\mcq{29}{If $A=\begin{bmatrix}3&1&3\\2&0&2\end{bmatrix}$ then order of $A^{t}$ is:}{3-by-2}{2-by-3}{3-by-3}{2-by-2}
\mcq{30}{$\begin{vmatrix}3&1\\0&4\end{vmatrix}=$}{11}{12}{13}{-11}
\mcq{31}{If $f(x)=\dfrac{5x-6}{3}$ then what is the value of $f(3)$?}{-1}{3}{9}{15}
\mcq{32}{A function f from X to Y represented by:}{$f:XY$}{$f:Y\rightarrow X$}{$f:X\rightarrow Y$}{$f:\dfrac{X}{Y}$}
\mcq{33}{$(fog)(x)=$}{$(f+g)(x)$}{$(f-g)(x)$}{$f(g(x))$}{$f(x)\div g(x)$}
\mcq{34}{If $f(x)=2x+3$, $g(x)=x+1$, then $f(x)+g(x)=$}{3x}{3x+4}{4}{$2x^{2}+3$}
\mcq{35}{If $f(x)=5x+2$, $h(x)=2x-2$, then $f(x)-h(x)=$}{3x}{$5x^{2}-4$}{3x+4}{-3x-4}
\mcq{36}{If $f(x)=3x+1$, $g(x)=2x$, then $g(x)\times f(x)=$}{6x+2x}{$5x^{2}+1$}{x+1}{$6x^{2}+2x$}
\mcq{37}{If $f(x)=x^{2}-4$, $g(x)=x+2$, $x\neq-2$, then $\dfrac{f(x)}{g(x)}=$}{$\dfrac{1}{x-2}$}{$\dfrac{1}{x+2}$}{x+2}{x-2}
\mcq{38}{What is the shape of an absolute value function:}{U-shaped}{V-shaped}{L-shaped}{M-shaped}
\mcq{39}{If a graph represent a function, then every vertical line must intersect it at:}{4 points}{3 points}{2 points}{1 point}
\mcq{40}{If $f(x)=x^{3}$, the $f(-2)=$}{-8}{8}{4}{-610}
\mcq{41}{The expression $\dfrac{2x-1}{x^{2}+4}$ is:}{A polynomial}{An algebraic fraction}{A numerical fraction}{An equation}
\mcq{42}{Lowest form of $\dfrac{18a^{2}bc^{2}}{12ac}$ is:}{$\dfrac{9abc}{6}$}{$\dfrac{3abc}{2}$}{$\dfrac{18b}{12}$}{$\dfrac{3}{2}$}
\mcq{43}{Lowest form of $\dfrac{15xyz}{10}$ is:}{$\dfrac{3xyz}{2}$}{$\dfrac{3}{2}$}{xyz}{$\dfrac{xyz}{2}$}
\mcq{44}{Simplified form of $\dfrac{3x+1}{5}+\dfrac{2x+3}{5}$ is}{$\dfrac{5x+4}{10}$}{$\dfrac{5x+4}{5}$}{$\dfrac{5x+4}{25}$}{5x+4}
\mcq{45}{Simplified form of $\left(\dfrac{x^{2}-y^{2}}{x+y}\right)(x-y)$ is:}{$(x+y)^{2}$}{$x^{2}+y^{2}$}{$x-y$}{$(x-y)^{2}$}
\mcq{46}{$(x^{3}-y^{3})\div(x-y)$ in simplified form is:}{$x^{2}-xy+y^{2}$}{$x^{2}+xy+y^{2}$}{$x^{2}-y^{2}$}{$x^{2}+y^{2}$}
\mcq{47}{An equation of the form $\dfrac{5}{x}+\dfrac{1-x}{3}=\dfrac{1}{6x}$ is called:}{Radical equation}{Reciprocal equation}{Fractional equation}{Exponential equation}
\mcq{48}{An equation of the form $5^{x}+64.5^{-x}-20=0$ is called:}{Radical equation}{Exponential equation}{Reciprocal equation}{Fractional equation}
\mcq{49}{Roots of $y^{2}-24y+128=0$ are:}{8,-16}{8,16}{-8,16}{-8,-16}
\mcq{50}{Linear factors of $3x^{2}+10x+3=0$ are:}{$(x+3),(3x+1)$}{$(x+3),(3x-1)$}{$(x-3),(3x+1)$}{$(x-3),(3x-1)$}
\mcq{51}{What is the value of $\cot60^{o}$:}{0}{1}{$\sqrt{3}$}{$\dfrac{1}{\sqrt{3}}$}
\mcq{52}{All trigonometric ratios are positive in:}{I-quadrant}{II-quadrant}{III-quadrant}{IV-quadrant}
\mcq{53}{$\mathrm{cosec}\,\theta$ is positive in:}{I \& III-quadrants}{II \& IV-quadrants}{I \& II-quadrants}{I \& IV quadrants}
\mcq{54}{The law of sines is:}{$\dfrac{b}{\sin\alpha}=\dfrac{a}{\sin\beta}=\dfrac{c}{\sin\gamma}$}{$\dfrac{a}{\sin\alpha}=\dfrac{c}{\sin\beta}=\dfrac{b}{\sin\gamma}$}{$\dfrac{a}{\sin\alpha}=\dfrac{b}{\sin\beta}=\dfrac{c}{\sin\gamma}$}{$\dfrac{a}{\sin\alpha}=\dfrac{a}{\sin\beta}=\dfrac{a}{\sin\ldots}$}
\mcq{55}{The law of cosines is:}{$\cos\alpha=\dfrac{b^{2}+a^{2}-c^{2}}{2ab}$}{$\cos\beta=\dfrac{c^{2}+b^{2}-a^{2}}{2cb}$}{$\cos\gamma=\dfrac{a^{2}+b^{2}-c^{2}}{2ab}$}{$\cos\gamma=\dfrac{ca^{2}+b^{2}-c}{2ab}$}
\mcq{56}{Area of $\Delta$ABC is:}{$\dfrac{1}{2}ab\sin\alpha\ \dfrac{1}{2}ab\sin\alpha$}{$\dfrac{1}{2}bc\sin\beta\ \dfrac{1}{2}bc\sin\beta$}{$\dfrac{1}{2}\sin\gamma\ \dfrac{1}{2}\sin\gamma$}{$\dfrac{1}{2}ac\sin\beta\ \dfrac{1}{2}ac\sin\beta$}
\mcq{57}{Bearing is measured from:}{East}{West}{North}{South}
\mcq{58}{Bearing is written as a:}{1 figure}{2 figures}{3 figures}{4 figures}
\mcq{59}{$\tan(180^{o}-\theta)=$}{$\tan\theta$}{$\cot\theta$}{$-\cot\theta$}{$-\tan\theta$}
\mcq{60}{$\sin(90^{o}+\theta)=$}{$\sin\theta$}{$-\sin\theta$}{$\cos\theta$}{$-\cos\theta$}
\mcq{61}{Distance of a point on the circumference to the centre of the circle is called:}{Radius}{Arc}{Chord}{Tangent}
\mcq{62}{Radii of same circles are:}{All unequal}{All equal}{Half of each chord}{Double of the diameter}
\mcq{63}{The boundary of the circle is called:}{Chord}{Segment}{Circumference}{Diameter}
\mcq{64}{Any part of a circumference is called:}{Chord}{Diameter}{Radius}{Arc}
\mcq{65}{A chord passing through the centre of the circle is called:}{Radius}{Diameter}{Secant}{Circumference}
\mcq{66}{In the given figure, what is AB (image):}{Diameter}{Tangent}{Chord}{Arc}
\mcq{67}{In the given figure, major arc is (image):}{$m\overset{\frown}{AB}$}{$m\overset{\frown}{BC}$}{$m\overset{\frown}{BDC}$}{$m\overset{\frown}{AD}$}
\mcq{68}{...... circle (s) can pass through the three non-collinear points:}{One}{Two}{Three}{Four}
\mcq{69}{Perpendicular bisector of a chord always passes through the ...... of circle:}{Arc}{Radius}{Centre}{Circumference}
\mcq{70}{The sum of the measures of central angles of a circle is:}{$90^{o}$}{$180^{o}$}{$270^{o}$}{$360^{o}$}
\mcq{71}{Tangent is a line that touches the circumference of the circle at:}{Many points}{Three points}{Two points}{One point}
\mcq{72}{...... tangents can be drawn to a circle from a point outside the circle:}{One}{Two}{Three}{Many}
\mcq{73}{....... tangents can be drawn to a circle from centre of the circle:}{Zero}{One}{Two}{Three}
\mcq{74}{In circle, radius and tangent are:}{Parallel}{Equal}{Perpendicular}{Zero}
\mcq{75}{An angle in a segment greater than a semicircle is ...... angle:}{An acute}{Right}{An obtuse}{Straight}
\mcq{76}{The angle subtended by the arc at the centre of the circle is called:}{Acute angle}{Central angle}{Right angle}{Complete angle}
\mcq{77}{What is the measure of an angle inscribed in a semi-circle?:}{$45^{o}$}{$60^{o}$}{$90^{o}$}{$120^{o}$}
\mcq{78}{Any two angles in the same segment of the circle are:}{Supplementary}{Complementary}{Equal}{Zero}
\mcq{79}{Area of the sector of the circle = .......:}{$\dfrac{1}{2}r\theta$}{$\dfrac{3}{2}r^{2}\theta$}{$\dfrac{1}{2}r^{2}\theta$}{$\dfrac{3}{2}r\theta$}
\mcq{80}{If r=6cm and $\theta=2$ radians, then arc length is:}{12 cm}{8 cm}{6 cm}{3 cm}
\mcq{81}{In the adjacent, $\overleftrightarrow{ABC}$ figure, is:}{A chord}{An arc}{A tangent}{A secant}
\mcq{82}{There are ...... types of arcs:}{2}{3}{4}{5}
\mcq{83}{Right bisector of the chord of a circle always passes through the:}{Diameter}{Non-collinear points}{Radius}{Centre}
\mcq{84}{A circle has only one:}{Chord}{Centre}{Diameter}{Secant}
\mcq{85}{The point where two tangents meet outside a circle forms:}{A semicircle}{A diameter}{A radius}{An angle}
\mcq{86}{Two equal tangents from a point to a circle can be drawn when the point is ...... the cirlce:}{On}{Inside}{Outside}{At centre of}
\mcq{87}{Tangents drawn from a single external point to a circle are:}{Unequal in length}{Perpendicular to each other}{Equal in length}{Inside the crcle}
\mcq{88}{At least how many chords are needed to locate the centre of the circle:}{1}{2}{3}{4}
\mcq{89}{To draw a tangent at middle point of an arc, first step is to:}{Join endpoints of an arc}{Draw radius}{Draw a line perpendicular to chord}{Draw an angle}
\mcq{90}{The angle between the radius and a tangent at the point of contact is:}{$30^{o}$}{$60^{o}$}{$90^{o}$}{$120^{o}$}
\mcq{91}{....... is used to get the comulative frequencies:}{Addition}{Square root}{Multiplication}{Division}
\mcq{92}{Second quartile represents:}{Mean}{Mode}{Median}{Variance}
\mcq{93}{First quartile divides the data into ...... equal parts:}{Two}{Three}{Four}{Ten}
\mcq{94}{Differece between the highest and the lowest values is called:}{Mean}{Variance}{Range}{Standard deviation}
\mcq{95}{Scatter diagram represents a relationship between ....... variables:}{Five}{Two}{Three}{Four}
\mcq{96}{...... is measure of dispersion:}{Mean}{Median}{Mode}{Variance}
\mcq{97}{Positve square root of variance is called:}{Mean}{Median}{Standard deviation}{Range}
\mcq{98}{...... is not a measure of dispersion:}{Range}{Arithmetic mean}{Variance}{Standard deviation}
\mcq{99}{Variance of the data 8,8,8,8,8,8 is:}{0}{16}{8}{48}
\mcq{100}{Range of first 20 natural numbers is:}{20}{10}{19}{30}
\mcq{101}{The probability of impossible event is:}{0}{1}{2}{-1}
\mcq{102}{What is the probability of getting a head when a fair coin is tossed once:}{0}{0.25}{0.5}{1}
\mcq{103}{The sum of probabilities of all possible outcomes of an experiment is:}{0.5}{0.25}{1}{0.4}
\mcq{104}{If P(A)=0.6 then the probability of event A not happening is:}{0.4}{0.6}{1}{1.6}
\mcq{105}{Two events are said to be mutually exclusive if:}{They can happed at the same time}{One affects the other}{They cannot happen together}{They are always equal}
\mcq{106}{If one coin is tossed and one dice is rolled, then the number of sample point are:}{3}{6}{2}{12}
\mcq{107}{What is the probability of sure event?}{0}{1}{$\dfrac{2}{3}$}{$\dfrac{4}{5}$}
\mcq{108}{The probability of getting a number greater than 6 on dice is:}{0.33}{1}{0}{0.5}
\mcq{109}{...... outcomes are possible when we draw a card from deck of cards:}{13}{1}{52}{26}
\mcq{110}{If P(E)=0.07 then the probability of Not E is:}{0.95}{0.89}{0.93}{0.90}

\end{document}
